\section{Appendix: Scope and Qualifications}\label{app:scope}

\subsection{Controlling Liturgical Books, Territory, and As-Of Date}

This seasonal handbook covers Advent in the Postconciliar Roman Rite as celebrated in the dioceses of the United States as of 27 November 2026. Its controlling liturgical authorities are:

\begin{itemize}
  \item \textit{Missale Romanum}, third typical edition (Vatican City, 2002; emended reprint, 2008), governed by the \textit{Institutio Generalis Missalis Romani} (GIRM) and the \textit{Normae Universales de Anno Liturgico et de Calendario} (1969, promulgated with Pope Paul~VI's motu proprio \textit{Mysterii Paschalis}).
  \item \textit{The Roman Missal}, Third Typical Edition (United States English, implemented 27 November 2011), together with the \textit{General Instruction of the Roman Missal} (U.S.~English edition, 2011, emended 2021).
  \item \textit{Ordo Lectionum Missae}, second typical edition (1981), and the United States \textit{Lectionary for Mass} (1998--2002), for Sundays (Years~A, B, and~C) and weekdays of Advent (nos.~175--200).
  \item \textit{Liturgia Horarum} (promulgated by Pope Paul~VI's apostolic constitution \textit{Laudis canticum}, 1 November 1970) and the United States \textit{Liturgy of the Hours}.
  \item The 1983 \textit{Code of Canon Law} (cc.~1246--1253), Pope Paul~VI's \textit{Paenitemini} (1966), the United States bishops' \textit{Pastoral Statement on Penance and Abstinence} (1966), their complementary norm on canons~1252--1253 (1983/1984), and the United States Proper Calendar (including December~8 Immaculate Conception and December~12 Our Lady of Guadalupe).
\end{itemize}

This handbook is a study and prayer companion, not an official liturgical book or a substitute for a diocesan ordo. Wherever a diocesan calendar, bishop's decree, or local parish schedule introduces a particular feast, transfer, or pastoral notice, follow your diocese and parish.

\subsection{Calendar Arithmetic and Multi-Year Verification (2020--2120)}

All calendar dates, cycle assignments, weekday counts, and overlap frequencies stated in this handbook were programmatically verified across 101 consecutive Advent seasons (2020 through 2120) and cross-checked on fifteen boundary dates against independent liturgical-calendar computations. Over that 101-year span:

\begin{itemize}
  \item The First Sunday of Advent falls between \textbf{November~27} (earliest; 15 seasons) and \textbf{December~3} (latest; 14 seasons).
  \item Advent lasts between \textbf{22 days} (shortest; 14 seasons, opening December~3) and \textbf{28 days} (longest; 15 seasons, opening November~27) from the First Sunday through the morning of December~24.
  \item Week~III weekdays occurring before December~17 range from \textbf{0} (28 seasons, whenever the Third Sunday falls on December~16 or~17) to \textbf{5} (15 seasons, whenever the First Sunday falls on November~27). Because Saturday of Week~III can never fall earlier than December~17, the Lectionary contains no Week~III Saturday.
  \item Privileged weekdays (December~17--24) occurring at Mass number \textbf{7} in 87 seasons (one displaced by the Fourth Sunday) and \textbf{6} in 14 seasons (when the Third Sunday falls on December~17 and the Fourth Sunday on December~24).
  \item December~8 falls on the Second Sunday of Advent \textbf{14 times} (2024, 2030, 2041, 2047, etc.), transferring the Solemnity of the Immaculate Conception to Monday, December~9.
  \item December~12 falls on the Third Sunday of Advent \textbf{15 times} (2021, 2027, 2032, 2038, 2049, etc.), displacing the United States Feast of Our Lady of Guadalupe without transfer.
\end{itemize}

\begin{fourstable}{Advent Year}{Cycle \& Dates}{Length / Wk~3}{United States Notes}
2025 & Yr~A: Nov~30, Dec~7, 14, 21 & 25 days / 2 Wk~3 & Dec~8 Mon (Oblig.); Dec~12 Fri \\
\addlinespace
\textbf{2026} & \textbf{Yr~B: Nov~29, Dec~6, 13, 20} & \textbf{26 days / 3 Wk~3} & \textbf{Dec~8 Tue (Oblig.); Dec~12 Sat; Dec~24 Thu} \\
\addlinespace
2027 & Yr~C: Nov~28, Dec~5, 12, 19 & 27 days / 4 Wk~3 & Dec~8 Wed (Oblig.); \textbf{Dec~12 Sun (Guadalupe omitted)} \\
\addlinespace
2028 & Yr~A: Dec~3, 10, 17, 24 & \textbf{22 days} / 0 Wk~3 & Shortest Advent; \textbf{Dec~24 Adv~4 Sun} (6 priv.~wkdays) \\
\addlinespace
2029 & Yr~B: Dec~2, 9, 16, 23 & 23 days / 0 Wk~3 & Dec~8 Sat (Oblig.); Dec~12 Wed \\
\addlinespace
2030 & Yr~C: Dec~1, 8, 15, 22 & 24 days / 1 Wk~3 & \textbf{Dec~8 Sun $\to$ Mon Dec~9} (no transferred oblig.) \\
\addlinespace
2033 & Yr~C: Nov~27, Dec~4, 11, 18 & \textbf{28 days} / 5 Wk~3 & Longest Advent; all 5 Wk~3 wkdays (Dec~12 Mon Guadalupe) \\
\end{fourstable}

\subsection{Bible Numbering and Copyright Boundaries}

Modern postconciliar liturgical books cite the Psalms and Old Testament according to the Hebrew / \textit{Nova Vulgata} numbering, whereas the public-domain Douay-Rheims Bible follows the Clementine Vulgate. Throughout this handbook, Psalm numbers are given in Hebrew numbering first with the Vulgate number in parentheses where they differ---for example, Psalm~85 (84) or Psalm~25 (24). Likewise, six Advent Old Testament passages carry shifted chapter or verse numbers between the modern Lectionary and Douay-Rheims: Isaiah~63:19b; 64:2b--7 (Vulgate~64:1, 3--8), Isaiah~9:1--6 (Vulgate~9:2--7), Micah~5:1--4a (Vulgate~5:2--5a), Malachi~3:1--4, 23--24 (Vulgate~3:1--4; 4:5--6), Zechariah~2:14--17 (Douay-Rheims~2:10--13), and Judith~13:18bcde, 19 (Douay-Rheims~13:23--25); and 1--2~Samuel appear as 1--2~Kings in Douay-Rheims.

No copyrighted English text from \textit{The Roman Missal} (ICEL, 2010/2011), the United States \textit{Lectionary for Mass} (Confraternity of Christian Doctrine / NAB), or the \textit{Liturgy of the Hours} (ICEL) is reproduced in this handbook. All English Scripture quotations come from the tracked public-domain Douay-Rheims text, and all complete vocal prayers, hymns, antiphons, and proclamations in Appendix~B come from tracked pre-1928 public-domain Catholic editions.

\subsection{Source Verification Bounds}

Primary historical, patristic, conciliar, papal, canonical, and devotional quotations in this handbook are verified against tracked primary editions recorded in the companion source-verification apparatus. Where current liturgical or canonical rules extend beyond individually excerpted editions---specifically the moderation of flowers, music, and vestment colors in GIRM~nos.~305, 313, and~346, the rank of the United States Feast of Our Lady of Guadalupe on December~12, and the post-2024 clarification by the Dicastery for Divine Worship and the Discipline of the Sacraments (Note of 23 January 2025, Prot.~n.~51/25, issued after consultation with the Dicastery for Legislative Texts) establishing that an occasional transfer of a holy day of obligation from Sunday does not transfer the Mass precept to the day \textit{ad quem}---those points are documented explicitly in the edition's source audit and remain subject to local diocesan publication.
